\section{分子光吸収・選択則と光異性化}

\subsection{平成30年度 問14}
\label{q:h30-reg-applied-14}

\noindent\textbf{原典：} sources/H30/17\_applied\_science.pdf、PDF物理5頁。

\textbf{公式正答：} 3（single）。

\par\textbf{検討未確定：} 視覚一次教科書の公開検索収録で光反応を確認。各分子全スペクトル・光合成全経路の本文確認は未完了。

\paragraph{問題の要旨} 分子光吸収の選択則・光異性化・線幅を比較する。

\Needspace{4\baselineskip}\paragraph{確認した条件・記述（独自要約）}\mbox{}\par

\noindent\textbf{判断対象：} 分子光吸収の選択則・光異性化・線幅を比較する。\par

\noindent\textbf{原典確認：} 全5候補・指数・否定・記号を保存原典22頁画像と照合。\par

\paragraph{問われる核心} 分子光吸収の選択則・光異性化・線幅を比較する。

\paragraph{必要知識} 視物質の11cisレチナールは可視光吸収後に異性化し信号へつながるので3。

\paragraph{考え方・正答を選ぶ理由（査読待ち初稿）} 視物質の11cisレチナールは可視光吸収後に異性化し信号へつながるので3。CO2の振動IR吸収、IRとRaman選択則、クロロフィルの励起とNADP+の電子受容、運動・衝突による線幅をそれぞれ別の物理機構として確認する。選択則は許容遷移を扱い、線幅が常に一定という保証ではない。

\Needspace{4\baselineskip}\paragraph{選択肢の検討（査読待ち初稿）}\mbox{}\par

\noindent\textbf{1：選ばない候補} CO2の振動は可視光を吸収して温室効果。温室効果の主要な振動吸収は赤外域。\par

\noindent\textbf{2：選ばない候補} IRは分極率変化で双極子変化ではない。IRは双極子変化、Ramanは分極率変化で取り違え。\par

\noindent\textbf{3：選ぶ候補} 11cisレチナールの光異性化で視覚。視物質中の11cis→alltransの光反応に対応。\par

\noindent\textbf{4：選ばない候補} クロロフィル又はNADP+が光子で一重項励起。クロロフィルの励起と電子受容体NADP+を同列の光捕集体にしない。\par

\noindent\textbf{5：選ばない候補} 分子運動や衝突で線幅は変わらない。Doppler幅・衝突幅等が存在する。\par

\paragraph{再利用する考え方} 軌道と電子状態、格子と基底、単位胞原子数と配位数、定義の分母を区別して候補を全て評価する。

\paragraph{共通理論} \hyperref[theory:atomic:optical]{第\ref*{theory:atomic:optical}節}

\par 参考：\cite{atomic_retinal,atomic_ir_raman,atomic_linewidth,atomic_photosynthesis,atomic_ir_co2}。

\par 難易度 2（校準初稿）、優先度未評価。

\section{電磁波領域・光子エネルギー}

\subsection{平成24年度 問3}
\label{q:h24-reg-applied-03}

\noindent\textbf{原典：} sources/H24/17\_applied\_science.pdf、PDF物理1頁。

\textbf{公式正答：} 1（single）。

\paragraph{問題の要旨} 赤外・可視・紫外・γ線の4領域を波長の長い順に並べる。

\Needspace{4\baselineskip}\paragraph{確認した条件・記述（独自要約）}\mbox{}\par

\noindent\textbf{方向：} 波長の長い順。\par

\noindent\textbf{対象：} 4つの電磁波領域。光子エネルギーを波の総エネルギーと区別。\par

\paragraph{問われる核心} 電磁波の波長・周波数・光子エネルギーの対応

\paragraph{必要知識} 電磁波の波長・周波数・光子エネルギーの対応。\(c=f\lambda,\quad E=hf=hc/\lambda\)より波長と光子エネルギーは逆順となる。波長の長い順は赤外→可視→紫外→γで1。γ線を周波数の小さい領域として扱わない。

\paragraph{考え方・正答を選ぶ理由（査読待ち初稿）} \(c=f\lambda,\quad E=hf=hc/\lambda\)より波長と光子エネルギーは逆順となる。波長の長い順は赤外→可視→紫外→γで1。γ線を周波数の小さい領域として扱わない。

\Needspace{4\baselineskip}\paragraph{選択肢の検討（査読待ち初稿）}\mbox{}\par

\noindent\textbf{1：最適} 赤外線→可視光線→紫外線→γ線。指定した順序は赤外→可視→紫外→γであり、この候補は整合する。\par

\noindent\textbf{2：不適切} 紫外線→可視光線→赤外線→γ線。指定した順序は赤外→可視→紫外→γであり、この候補は少なくとも隣接2領域の大小関係が逆又は途中に不適切な領域を挟む。\par

\noindent\textbf{3：不適切} 赤外線→γ線→紫外線→可視光線。指定した順序は赤外→可視→紫外→γであり、この候補は少なくとも隣接2領域の大小関係が逆又は途中に不適切な領域を挟む。\par

\noindent\textbf{4：不適切} γ線→紫外線→可視光線→赤外線。指定した順序は赤外→可視→紫外→γであり、この候補は少なくとも隣接2領域の大小関係が逆又は途中に不適切な領域を挟む。\par

\noindent\textbf{5：不適切} γ線→赤外線→紫外線→可視光線。指定した順序は赤外→可視→紫外→γであり、この候補は少なくとも隣接2領域の大小関係が逆又は途中に不適切な領域を挟む。\par

\paragraph{再利用する考え方} 電磁波の波長・周波数・光子エネルギーの対応を暗記だけでなく、位相・幾何・符号と成立条件から導く。選択肢番号を類似年度へそのまま移さない。

\paragraph{共通理論} \hyperref[theory:optics:spectrum]{第\ref*{theory:optics:spectrum}節}

\par 参考：\cite{os_spectrum}。

\paragraph{関連問題} \hyperref[q:r02-reg-applied-07]{R02-REG-APPLIED-07}

\par 難易度 2（校準初稿）、優先度未評価。

\subsection{令和2年度 問7}
\label{q:r02-reg-applied-07}

\noindent\textbf{原典：} sources/R02/17\_applied\_science.pdf、PDF物理3頁。

\textbf{公式正答：} 4（single）。

\paragraph{問題の要旨} 赤外・可視・紫外・γ線の4領域を光子エネルギーの大きい順に並べる。

\Needspace{4\baselineskip}\paragraph{確認した条件・記述（独自要約）}\mbox{}\par

\noindent\textbf{方向：} 光子エネルギーの大きい順。\par

\noindent\textbf{対象：} 4つの電磁波領域。光子エネルギーを波の総エネルギーと区別。\par

\paragraph{問われる核心} 電磁波の波長・周波数・光子エネルギーの対応

\paragraph{必要知識} 電磁波の波長・周波数・光子エネルギーの対応。\(c=f\lambda,\quad E=hf=hc/\lambda\)より波長と光子エネルギーは逆順となる。光子エネルギーの大きい順はγ→紫外→可視→赤外で4。γ線を周波数の小さい領域として扱わない。

\paragraph{考え方・正答を選ぶ理由（査読待ち初稿）} \(c=f\lambda,\quad E=hf=hc/\lambda\)より波長と光子エネルギーは逆順となる。光子エネルギーの大きい順はγ→紫外→可視→赤外で4。γ線を周波数の小さい領域として扱わない。

\Needspace{4\baselineskip}\paragraph{選択肢の検討（査読待ち初稿）}\mbox{}\par

\noindent\textbf{1：不適切} 赤外線→可視光線→紫外線→γ線。指定した順序はγ→紫外→可視→赤外であり、この候補は少なくとも隣接2領域の大小関係が逆又は途中に不適切な領域を挟む。\par

\noindent\textbf{2：不適切} 紫外線→可視光線→赤外線→γ線。指定した順序はγ→紫外→可視→赤外であり、この候補は少なくとも隣接2領域の大小関係が逆又は途中に不適切な領域を挟む。\par

\noindent\textbf{3：不適切} 赤外線→γ線→紫外線→可視光線。指定した順序はγ→紫外→可視→赤外であり、この候補は少なくとも隣接2領域の大小関係が逆又は途中に不適切な領域を挟む。\par

\noindent\textbf{4：最適} γ線→紫外線→可視光線→赤外線。指定した順序はγ→紫外→可視→赤外であり、この候補は整合する。\par

\noindent\textbf{5：不適切} γ線→赤外線→紫外線→可視光線。指定した順序はγ→紫外→可視→赤外であり、この候補は少なくとも隣接2領域の大小関係が逆又は途中に不適切な領域を挟む。\par

\paragraph{再利用する考え方} 電磁波の波長・周波数・光子エネルギーの対応を暗記だけでなく、位相・幾何・符号と成立条件から導く。選択肢番号を類似年度へそのまま移さない。

\paragraph{共通理論} \hyperref[theory:optics:spectrum]{第\ref*{theory:optics:spectrum}節}

\par 参考：\cite{os_spectrum}。

\paragraph{関連問題} \hyperref[q:h24-reg-applied-03]{H24-REG-APPLIED-03}

\par 難易度 2（校準初稿）、優先度未評価。

\section{光の干渉・回折格子}

\subsection{平成27年度 問3}
\label{q:h27-reg-applied-03}

\noindent\textbf{原典：} sources/H27/17\_applied\_science.pdf、PDF物理1頁。

\textbf{公式正答：} 1（single）。

\par\textbf{検討未確定：} 公式・独立解答1一致。複スリットN>2では主極大間に副極大・N−1個の暗線があり、全暗線の隣接間隔を主極大間隔と同一にできない。原典の括弧書きの範囲をAstraで確認。

\paragraph{問題の要旨} 複スリットによる干渉縞の間隔を広げる変更を選ぶ。

\Needspace{4\baselineskip}\paragraph{確認した条件・記述（独自要約）}\mbox{}\par

\noindent\textbf{配置：} スリット間隔d、本数N、強度I、波長λ、スクリーン距離L。\par

\noindent\textbf{方向：} 明線（もしくは暗線）の相互間隔を広げる変更。\par

\paragraph{問われる核心} 干渉主極大の位置と複スリット本数の役割

\paragraph{必要知識} 干渉主極大の位置と複スリット本数の役割。\(d\sin\theta=m\lambda\)、遠方・小角で\(\Delta x\simeq L\lambda/d\)より1。原典の「明線（もしくは暗線）」を複スリットの全ての暗線の個数まで含めた一様な間隔と読むと本数Nの影響もあり、主極大同士の間隔との区別を残す。

\paragraph{考え方・正答を選ぶ理由（査読待ち初稿）} \(d\sin\theta=m\lambda\)、遠方・小角で\(\Delta x\simeq L\lambda/d\)より1。原典の「明線（もしくは暗線）」を複スリットの全ての暗線の個数まで含めた一様な間隔と読むと本数Nの影響もあり、主極大同士の間隔との区別を残す。

\Needspace{4\baselineskip}\paragraph{選択肢の検討（査読待ち初稿）}\mbox{}\par

\noindent\textbf{1：最適} スリット間隔dを狭くする。。主極大間隔Lλ/dはdを小さくすると広がる。\par

\noindent\textbf{2：不適切} スリットの本数Nを増やす。。本数は主極大を狭く強くし、その間に副極大・暗線を増やすが、主極大同士の角間隔は変えない。\par

\noindent\textbf{3：不適切} スクリーンまでの距離Lを短くする。。Lを短くするとスクリーン上の主極大間隔は狭くなる。\par

\noindent\textbf{4：不適切} 入射光の強度Iを強くする。。線形領域では強度は明るさを変え、位相の極大位置を変えない。\par

\noindent\textbf{5：不適切} 単色光の波長λを短くする。。波長を短くすると主極大間隔は狭くなる。\par

\paragraph{再利用する考え方} 干渉主極大の位置と複スリット本数の役割を暗記だけでなく、位相・幾何・符号と成立条件から導く。選択肢番号を類似年度へそのまま移さない。

\paragraph{共通理論} \hyperref[theory:optics:grating]{第\ref*{theory:optics:grating}節}

\par 参考：\cite{os_grating,os_double}。

\par 難易度 2（校準初稿）、優先度未評価。

\subsection{令和元年度再試験 問3}
\label{q:r01-ret-applied-03}

\noindent\textbf{原典：} sources/R01\_retest/17\_applied\_science.pdf、PDF物理2頁。

\textbf{公式正答：} 5（single）。

\paragraph{問題の要旨} 反射位相、ニュートンリング、スリット、色収差、光ファイバの屈折率を判定する。

\Needspace{4\baselineskip}\paragraph{確認した条件・記述（独自要約）}\mbox{}\par

\noindent\textbf{方向：} 最も不適切。\par

\noindent\textbf{範囲：} 透明ガラス・単色干渉・通常の凸レンズ・光ファイバ。入射角等の細かな指定なし。\par

\paragraph{問われる核心} 屈折率の大小と位相・全反射・色収差

\paragraph{必要知識} 屈折率の大小と位相・全反射・色収差。コアからクラッドへの全反射条件は\(n_{\rm core}>n_{\rm clad}\)である。逆にした5が不適切。その他は通常の透明誘電体・近軸・正常分散という基礎光学モデルの範囲で適切であり、未指定の角度を問題の所与条件として追加しない。

\paragraph{考え方・正答を選ぶ理由（査読待ち初稿）} コアからクラッドへの全反射条件は\(n_{\rm core}>n_{\rm clad}\)である。逆にした5が不適切。その他は通常の透明誘電体・近軸・正常分散という基礎光学モデルの範囲で適切であり、未指定の角度を問題の所与条件として追加しない。

\Needspace{4\baselineskip}\paragraph{選択肢の検討（査読待ち初稿）}\mbox{}\par

\noindent\textbf{1：適切（通常のモデル）} 空気中のガラスの表面反射は位相反転し、裏面での内部反射は反転しない。。損失のない通常の誘電体のほぼ垂直な入射では、低屈折率から高屈折率の反射でπ反転し、逆方向では反転しない。原典に入射角指定はなく、全斜入射・偏光まで同じとする規則へ拡張しない。\par

\noindent\textbf{2：適切} ニュートンリング半径は照射波長が長いほど大きい。。同じ曲率半径・次数で半径の二乗が波長に比例する。反射明環と暗環の定数は違うが大小の方向は一致。\par

\noindent\textbf{3：適切} 2スリットの間隔を狭くすると遠方の干渉縞間隔は広がる。。遠方・小角でLλ/dであるから、dの減少で間隔が増える。\par

\noindent\textbf{4：適切（正常分散）} 凸レンズの平行光の焦点は紫より赤が遠い。。通常の分散では紫の屈折率が赤より大きく、同じ形状の凸レンズの焦点距離は紫が短い。異常分散を無制限に排除する一般則ではない。\par

\noindent\textbf{5：不適切} 光ファイバはクラッドの屈折率をコアより大きくして全反射を利用する。。全反射には高屈折率側から低屈折率側への入射が必要。コアの方が高く、記述が逆。\par

\paragraph{再利用する考え方} 屈折率の大小と位相・全反射・色収差を暗記だけでなく、位相・幾何・符号と成立条件から導く。選択肢番号を類似年度へそのまま移さない。

\paragraph{共通理論} \hyperref[theory:optics:reflection]{第\ref*{theory:optics:reflection}節}, \hyperref[theory:optics:newton]{第\ref*{theory:optics:newton}節}, \hyperref[theory:optics:grating]{第\ref*{theory:optics:grating}節}, \hyperref[theory:optics:refraction]{第\ref*{theory:optics:refraction}節}

\par 参考：\cite{os_fiber,os_thinfilm,os_lens}。

\par 難易度 2（校準初稿）、優先度未評価。

\subsection{令和2年度 問5}
\label{q:r02-reg-applied-05}

\noindent\textbf{原典：} sources/R02/17\_applied\_science.pdf、PDF物理2頁。

\textbf{公式正答：} 5（single）。

\paragraph{問題の要旨} 単色光の第一回折明線の角度から回折格子1 mm当たりの本数を求める。

\Needspace{4\baselineskip}\paragraph{確認した条件・記述（独自要約）}\mbox{}\par

\noindent\textbf{光：} 格子へ垂直に入射、波長650 nm。\par

\noindent\textbf{回折：} 第一明線30度。\par

\noindent\textbf{方向：} 1 mm当たりの格子本数に最も近い値。\par

\paragraph{問われる核心} 回折格子の光路差と格子定数の逆数

\paragraph{必要知識} 回折格子の光路差と格子定数の逆数。\(d\sin30^\circ=650\ {\rm nm}\)からd=\allowbreak{}1300 nm=\allowbreak{}0.0013 mm。したがって\(1/d\simeq769.2\) 本/mmとなり5。格子間隔と本数を同じ量としない。

\paragraph{考え方・正答を選ぶ理由（査読待ち初稿）} \(d\sin30^\circ=650\ {\rm nm}\)からd=\allowbreak{}1300 nm=\allowbreak{}0.0013 mm。したがって\(1/d\simeq769.2\) 本/mmとなり5。格子間隔と本数を同じ量としない。

\Needspace{4\baselineskip}\paragraph{選択肢の検討（査読待ち初稿）}\mbox{}\par

\noindent\textbf{1：不適切} 33 本/mm。格子間隔の逆数と第一次数の条件を満たさず、正しい約769本/mmから遠い。\par

\noindent\textbf{2：不適切} 65 本/mm。nmからmmへの変換と、間隔の逆数を確認する。正しい本数とは一桁以上違う。\par

\noindent\textbf{3：不適切} 130 本/mm。間隔1300 nmを単なる整数の本数へ読み替えていないか確認する。正しい約769本/mmではない。\par

\noindent\textbf{4：不適切} 390 本/mm。sin30°の因子を逆に扱う等の誤りで、本数がほぼ半分になっている。\par

\noindent\textbf{5：最適} 770 本/mm。d=\allowbreak{}1.3 µmに対して1 mm/d≈769本/mmで最も近い。\par

\paragraph{再利用する考え方} 回折格子の光路差と格子定数の逆数を暗記だけでなく、位相・幾何・符号と成立条件から導く。選択肢番号を類似年度へそのまま移さない。

\paragraph{共通理論} \hyperref[theory:optics:grating]{第\ref*{theory:optics:grating}節}

\par 参考：\cite{os_grating}。

\par 難易度 2（校準初稿）、優先度未評価。

\section{薄レンズの実像・虚像と符号}

\subsection{平成23年度 問7}
\label{q:h23-reg-applied-07}

\noindent\textbf{原典：} sources/H23/17\_applied\_science.pdf、PDF物理3頁。

\textbf{公式正答：} 2（single）。

\paragraph{問題の要旨} 凸レンズの物体位置、像の実虚と向き、物体移動と像距離の変化を表の5空欄へ対応させる。

\Needspace{4\baselineskip}\paragraph{確認した条件・記述（独自要約）}\mbox{}\par

\noindent\textbf{候補表：} 各行はア・イ・ウ・エ・オの順。\par

\noindent\textbf{条件：} 凸レンズ、焦点距離の内外と有限の物体位置。\par

\noindent\textbf{方向：} 表の適切な組合せ。\par

\paragraph{問われる核心} 薄レンズの符号と像距離の単調性

\paragraph{必要知識} 薄レンズの符号と像距離の単調性。\(b=fa/(a-f)\)より、a>fならb>0の倒立実像で、\(db/da=-f^2/(a-f)^2<0\)なので物体がレンズに近付くと実像は遠ざかる。0<a<fならb<0、倍率−b/a>0の正立虚像である。5空欄が一致するのは2。

\paragraph{考え方・正答を選ぶ理由（査読待ち初稿）} \(b=fa/(a-f)\)より、a>fならb>0の倒立実像で、\(db/da=-f^2/(a-f)^2<0\)なので物体がレンズに近付くと実像は遠ざかる。0<a<fならb<0、倍率−b/a>0の正立虚像である。5空欄が一致するのは2。

\Needspace{4\baselineskip}\paragraph{選択肢の検討（査読待ち初稿）}\mbox{}\par

\noindent\textbf{1：不適切} 外側／虚像／同じ／遠い／実像。焦点外の像を虚像とした点と正立とした点が誤り。内側では実像ではない。\par

\noindent\textbf{2：最適} 外側／実像／反対／近い／虚像。焦点外は倒立実像で、物体を焦点へ近づけると像は遠ざかる。内側は正立虚像であり、全空欄が整合。\par

\noindent\textbf{3：不適切} 内側／虚像／反対／遠い／実像。実像形成の対象位置と虚像の向きを取り違え、内側の像を実像としている。\par

\noindent\textbf{4：不適切} 内側／実像／同じ／遠い／虚像。焦点内の像は正立だが虚像である。最初の実像条件と像の向き、最後の実像が不整合。\par

\noindent\textbf{5：不適切} 焦点上／実像／反対／近い／虚像。物体が焦点上なら出射光は平行で、有限位置へ像はできない。最初の空欄が成立しない。\par

\paragraph{再利用する考え方} 薄レンズの符号と像距離の単調性を暗記だけでなく、位相・幾何・符号と成立条件から導く。選択肢番号を類似年度へそのまま移さない。

\paragraph{共通理論} \hyperref[theory:optics:lens]{第\ref*{theory:optics:lens}節}

\par 参考：\cite{os_lens}。

\paragraph{関連問題} \hyperref[q:h26-reg-applied-07]{H26-REG-APPLIED-07}

\par 難易度 2（校準初稿）、優先度未評価。

\subsection{平成26年度 問7}
\label{q:h26-reg-applied-07}

\noindent\textbf{原典：} sources/H26/17\_applied\_science.pdf、PDF物理3頁。

\textbf{公式正答：} 2（single）。

\paragraph{問題の要旨} 凸レンズの物体位置、像の実虚と向き、物体移動と像距離の変化を表の5空欄へ対応させる。

\Needspace{4\baselineskip}\paragraph{確認した条件・記述（独自要約）}\mbox{}\par

\noindent\textbf{候補表：} 各行はア・イ・ウ・エ・オの順。\par

\noindent\textbf{条件：} 凸レンズ、焦点距離の内外と有限の物体位置。\par

\noindent\textbf{方向：} 表の適切な組合せ。\par

\paragraph{問われる核心} 薄レンズの符号と像距離の単調性

\paragraph{必要知識} 薄レンズの符号と像距離の単調性。\(b=fa/(a-f)\)より、a>fならb>0の倒立実像で、\(db/da=-f^2/(a-f)^2<0\)なので物体がレンズに近付くと実像は遠ざかる。0<a<fならb<0、倍率−b/a>0の正立虚像である。5空欄が一致するのは2。

\paragraph{考え方・正答を選ぶ理由（査読待ち初稿）} \(b=fa/(a-f)\)より、a>fならb>0の倒立実像で、\(db/da=-f^2/(a-f)^2<0\)なので物体がレンズに近付くと実像は遠ざかる。0<a<fならb<0、倍率−b/a>0の正立虚像である。5空欄が一致するのは2。

\Needspace{4\baselineskip}\paragraph{選択肢の検討（査読待ち初稿）}\mbox{}\par

\noindent\textbf{1：不適切} 外側／虚像／同じ／遠い／実像。焦点外の像を虚像とした点と正立とした点が誤り。内側では実像ではない。\par

\noindent\textbf{2：最適} 外側／実像／反対／近い／虚像。焦点外は倒立実像で、物体を焦点へ近づけると像は遠ざかる。内側は正立虚像であり、全空欄が整合。\par

\noindent\textbf{3：不適切} 内側／虚像／反対／遠い／実像。実像形成の対象位置と虚像の向きを取り違え、内側の像を実像としている。\par

\noindent\textbf{4：不適切} 内側／実像／同じ／遠い／虚像。焦点内の像は正立だが虚像である。最初の実像条件と像の向き、最後の実像が不整合。\par

\noindent\textbf{5：不適切} 焦点上／実像／反対／近い／虚像。物体が焦点上なら出射光は平行で、有限位置へ像はできない。最初の空欄が成立しない。\par

\paragraph{再利用する考え方} 薄レンズの符号と像距離の単調性を暗記だけでなく、位相・幾何・符号と成立条件から導く。選択肢番号を類似年度へそのまま移さない。

\paragraph{共通理論} \hyperref[theory:optics:lens]{第\ref*{theory:optics:lens}節}

\par 参考：\cite{os_lens}。

\paragraph{関連問題} \hyperref[q:h23-reg-applied-07]{H23-REG-APPLIED-07}

\par 難易度 2（校準初稿）、優先度未評価。

\subsection{平成28年度 問5}
\label{q:h28-reg-applied-05}

\noindent\textbf{原典：} sources/H28/17\_applied\_science.pdf、PDF物理2頁。

\textbf{公式正答：} 3（single）。

\paragraph{問題の要旨} 2枚の凸レンズの中間虚像を次の物体として、第二レンズの焦点距離を求める。

\Needspace{4\baselineskip}\paragraph{確認した条件・記述（独自要約）}\mbox{}\par

\noindent\textbf{第一レンズ：} 焦点距離15 cm、物体は手前6 cm。\par

\noindent\textbf{第二レンズ：} 第一より20 cm後方、最終像は第二より後方60 cm。\par

\noindent\textbf{方向：} 第二レンズの焦点距離。\par

\paragraph{問われる核心} 2枚の薄レンズと中間像の符号・距離基準

\paragraph{必要知識} 2枚の薄レンズと中間像の符号・距離基準。第一レンズは\(1/15=1/6+1/b_1\)より\(b_1=-10\) cmの虚像。第二レンズは第一より20 cm右なので、この像は第二から30 cm左の実物体となる。最終像60 cmを使い\(f_2=(1/30+1/60)^{-1}=20\) cmで3。

\paragraph{考え方・正答を選ぶ理由（査読待ち初稿）} 第一レンズは\(1/15=1/6+1/b_1\)より\(b_1=-10\) cmの虚像。第二レンズは第一より20 cm右なので、この像は第二から30 cm左の実物体となる。最終像60 cmを使い\(f_2=(1/30+1/60)^{-1}=20\) cmで3。

\Needspace{4\baselineskip}\paragraph{選択肢の検討（査読待ち初稿）}\mbox{}\par

\noindent\textbf{1：不適切} 9 cm。30 cmと60 cmに対し1/f=\allowbreak{}1/30+1/60が成立しない。\par

\noindent\textbf{2：不適切} 18 cm。中間像を第一レンズ後方にある実像とみなすと物体距離を誤る。正しい値20 cmとは一致しない。\par

\noindent\textbf{3：最適} 20 cm。第二レンズから見た物体距離30 cm、実像60 cmを満たす。\par

\noindent\textbf{4：不適切} 21 cm。正しい物体・像距離をレンズ式へ入れた20 cmと一致しない。\par

\noindent\textbf{5：不適切} 26 cm。第一レンズの物体距離6 cmをそのまま第二レンズへ使う等の距離基準の混同ではないか確認する。正しい20 cmとは一致しない。\par

\paragraph{再利用する考え方} 2枚の薄レンズと中間像の符号・距離基準を暗記だけでなく、位相・幾何・符号と成立条件から導く。選択肢番号を類似年度へそのまま移さない。

\paragraph{共通理論} \hyperref[theory:optics:lens]{第\ref*{theory:optics:lens}節}, \hyperref[theory:optics:cascade]{第\ref*{theory:optics:cascade}節}

\par 参考：\cite{os_lens}。

\paragraph{関連問題} \hyperref[q:r06-reg-applied-05]{R06-REG-APPLIED-05}

\par 難易度 2（校準初稿）、優先度未評価。

\subsection{令和6年度 問5}
\label{q:r06-reg-applied-05}

\noindent\textbf{原典：} sources/R06/17\_applied\_science.pdf、PDF物理2頁。

\textbf{公式正答：} 5（single）。

\paragraph{問題の要旨} 焦点内の物体を凸レンズで見たときの像の位置・向き・実虚を求める。

\Needspace{4\baselineskip}\paragraph{確認した条件・記述（独自要約）}\mbox{}\par

\noindent\textbf{レンズ：} 焦点距離15 cmの凸レンズ。\par

\noindent\textbf{物体：} レンズ手前10 cm。\par

\noindent\textbf{方向：} 像位置と実虚・向きの組合せ。\par

\paragraph{問われる核心} 薄レンズ焦点内の正立虚像と符号

\paragraph{必要知識} 薄レンズ焦点内の正立虚像と符号。\(1/15=1/10+1/b\)よりb=−30 cm。像は入射側にあり、倍率\(M=-b/a=3>0\)なので正立虚像で5。負の像距離を負の物理的長さとせず、位置の符号として読む。

\paragraph{考え方・正答を選ぶ理由（査読待ち初稿）} \(1/15=1/10+1/b\)よりb=−30 cm。像は入射側にあり、倍率\(M=-b/a=3>0\)なので正立虚像で5。負の像距離を負の物理的長さとせず、位置の符号として読む。

\Needspace{4\baselineskip}\paragraph{選択肢の検討（査読待ち初稿）}\mbox{}\par

\noindent\textbf{1：不適切} 手前15 cmに正立虚像。。正立虚像という性質は合うが、像距離の大きさは30 cm。\par

\noindent\textbf{2：不適切} 手前15 cmに倒立実像。。像の距離、倒立、実像の全てが焦点内物体の像と異なる。\par

\noindent\textbf{3：不適切} 後方6 cmに倒立虚像。。光線の延長は物体側へ交わるので、位置・向き・6 cmの距離が不整合。\par

\noindent\textbf{4：不適切} 後方30 cmに倒立実像。。30 cmという大きさのみ合い、物体側という位置と正立虚像という性質が逆。\par

\noindent\textbf{5：最適} 手前30 cmに正立虚像。。b=-30 cm、倍率−b/a=\allowbreak{}3なので物体側30 cmの正立虚像。\par

\paragraph{再利用する考え方} 薄レンズ焦点内の正立虚像と符号を暗記だけでなく、位相・幾何・符号と成立条件から導く。選択肢番号を類似年度へそのまま移さない。

\paragraph{共通理論} \hyperref[theory:optics:lens]{第\ref*{theory:optics:lens}節}

\par 参考：\cite{os_lens}。

\paragraph{関連問題} \hyperref[q:h28-reg-applied-05]{H28-REG-APPLIED-05}

\par 難易度 2（校準初稿）、優先度未評価。

\section{薄膜反射位相・ニュートンリング}

\subsection{令和3年度 問5}
\label{q:r03-reg-applied-05}

\noindent\textbf{原典：} sources/R03/17\_applied\_science.pdf、PDF物理2頁。

\textbf{公式正答：} 3（single）。

\par\textbf{検討未確定：} 独立3・公式3一致。ただし3はm=\allowbreak{}0始まりの干渉次数で成立。m番目を1始まりの序数と読むならm−1/2が必要で候補なし。開始値を原典所与として補わず部分校正。Astraで反射/透過と位相反転も確認。

\paragraph{問題の要旨} 空気層の反射ニュートンリングで明環半径と曲率半径・波長・次数を対応させる。

\Needspace{4\baselineskip}\paragraph{確認した条件・記述（独自要約）}\mbox{}\par

\noindent\textbf{配置：} 曲率半径Rの平凸レンズを凸面下に平面ガラスへ置く。接点Oから距離x。\par

\noindent\textbf{光：} 上から垂直な単色光、波長λ。反射で同心円状の縞を観察。\par

\noindent\textbf{方向：} 原文はOからm番目の明環。次数の開始値は明示なし。\par

\paragraph{問われる核心} 反射位相・空気層の厚さ・明環の次数定義

\paragraph{必要知識} 反射位相・空気層の厚さ・明環の次数定義。球面の幾何から\(t=R-\sqrt{R^2-x^2}\simeq x^2/(2R)\)。空気層の一方の反射にπの相対位相があるため反射明環は\(2t=(n+1/2)\lambda\)、n=\allowbreak{}0,1,…で、式は3に対応する。原典の「m番目」を1からの序数と読むなら\(x^2/R=(m-1/2)\lambda\)となり、該当肢がない。公式3を保持し、次数記号の解釈を確定済みとしない。

\paragraph{考え方・正答を選ぶ理由（査読待ち初稿）} 球面の幾何から\(t=R-\sqrt{R^2-x^2}\simeq x^2/(2R)\)。空気層の一方の反射にπの相対位相があるため反射明環は\(2t=(n+1/2)\lambda\)、n=\allowbreak{}0,1,…で、式は3に対応する。原典の「m番目」を1からの序数と読むなら\(x^2/R=(m-1/2)\lambda\)となり、該当肢がない。公式3を保持し、次数記号の解釈を確定済みとしない。

\Needspace{4\baselineskip}\paragraph{選択肢の検討（査読待ち初稿）}\mbox{}\par

\noindent\textbf{1：不適切} \(2x^2/R=m\lambda\)。曲面の幾何による往復空気層の厚さはx\textsuperscript{2}/Rで、2x\textsuperscript{2}/Rではない。整数倍も反射明環に対応しない。\par

\noindent\textbf{2：不適切} \(2x^2/R=(m+1/2)\lambda\)。半整数は反射明環の条件だが、幾何の係数が2倍。\par

\noindent\textbf{3：最適} \(x^2/R=(m+1/2)\lambda\)。mを0,1,2,…とする干渉次数なら正しい。ただし原典のm番目の数え始めは明示されない。\par

\noindent\textbf{4：不適切} \(x^2/R=m\lambda\)。片方の反射のみ位相反転するため、整数倍は反射暗環の条件となる。\par

\noindent\textbf{5：不適切} \(x^2/(2R)=(m+1/2)\lambda\)。x\textsuperscript{2}/(2R)は片道の厚さtであり、明環の位相条件には往復2tを使う。\par

\paragraph{再利用する考え方} 反射位相・空気層の厚さ・明環の次数定義を暗記だけでなく、位相・幾何・符号と成立条件から導く。選択肢番号を類似年度へそのまま移さない。

\paragraph{共通理論} \hyperref[theory:optics:newton]{第\ref*{theory:optics:newton}節}, \hyperref[theory:optics:reflection]{第\ref*{theory:optics:reflection}節}

\par 参考：\cite{os_thinfilm}。

\par 難易度 2（校準初稿）、優先度未評価。
